\input{../tex-inputs/latex-leseansicht-vorspann}

\section[Gustav Schwarzkopf an Arthur Schnitzler, {[}15. 5. 1892?{]}]{L04278 Gustav Schwarzkopf an Arthur Schnitzler, {[}15. 5. 1892?{]}}
\nopagebreak\mylabel{L04278v}
\rehead{ }\normalsize\beginnumbering\briefempfaengerindex{Schnitzler, Arthur@\textsc{Schnitzler, Arthur}!zzzSchwarzkopf, Gustav@\emph{von Gustav Schwarzkopf}!1892-05-151@{{[}15. 5. 1892?{]}}|(be}
\toendnotes[C]{\smallbreak\pagebreak[2]}
\correspDesc{Versand  durch Gustav Schwarzkopf am [15. 5. 1892?] in Wien
\newline{}Erhalt  durch Arthur Schnitzler am [15. 5. 1892?] in Wien}\toendnotes[C]{\smallbreak}
\Standort{CUL, Schnitzler, B 96a.}
\physDesc{Visitenkarte, 132 Zeichen
\newline{}Handschrift: schwarze Tinte, deutsche Kurrent}\toendnotes[C]{\smallbreak}
\pstart
           \centering{}{\pb}\textcolor{gray}{\textbf{Gustav Schwarzkopf}}\pend
           
\pstart{}{\pb}Sehr geehrter Herr
                  Doctor!\pend\vspace{0.5em}
\pstart
           \label{K_L04278-1v}\edtext{Minden\orgindex{Heinrich Minden@Heinrich Minden|pw} iſt eigenſi{\geminationn}ig und will nun einmal nicht}{\lemma{\textnormal{\emph{Minden … nicht}}}\Cendnote{\textnormal{Diese Visitenkarte von Schwarzkopf\pwindex{Schwarzkopf, Gustav 7.\,11.\,1853 Wien – 13.\,11.\,1939 ebd.@\textsc{Schwarzkopf, Gustav} (7.\,11.\,1853 Wien – 13.\,11.\,1939 ebd.), \emph{Schriftsteller}|pwk} dürfte dem 
               von \emph{Heinrich Minden}\orgindex{Heinrich Minden@Heinrich Minden|pwk} zurückgesandten Manuskript von \emph{Anatol}\pwindex{Schnitzler, Arthur 15.\,5.\,1862 Wien – 21.\,10.\,1931 ebd.@\textsc{Schnitzler, Arthur} (15.\,5.\,1862 Wien – 21.\,10.\,1931 ebd.), \emph{Schriftsteller, Mediziner}!Anatol@\strich\emph{Anatol}|pwk} beigelegen sein, vgl. A. S.: \emph{Tagebuch}, 15. 5. 1892.
            }}}\label{K_L04278-1}. Mit aufrichtigem Bedauern und den beſten Grüßen\pend
           
\pstart
           Ihr{\\[\baselineskip]}ergebener{\\[\baselineskip]}\spacefill\mbox{S}\pend
           \leftskip=0em{}\selectlanguage{ngerman}\endnumbering\briefempfaengerindex{Schnitzler, Arthur@\textsc{Schnitzler, Arthur}!zzzSchwarzkopf, Gustav@\emph{von Gustav Schwarzkopf}!1892-05-151@{{[}15. 5. 1892?{]}}|)be}\mylabel{L04278h}
\begin{anhang}
\end{anhang}\newcommand{\dateiname}{L04278}\newcommand{\titel}{Gustav Schwarzkopf an Arthur Schnitzler, [15. 5. 1892?]}\newcommand{\editorInnen}{Herausgegeben von Jahnke, SelmaMüller, Martin Anton}\input{../tex-inputs/latex-leseansicht-abspann}
